\section{Introduction}\label{sec:intro}

Deductive verifiers fall roughly
into two camps.
%
Satisfiability Modulo Theory (SMT)
based verifiers (\eg \dafny and \fstar)
use fast decision procedures to
automate the verification of programs that
only require reasoning over a fixed
set of theories like linear arithmetic,
string, set and bitvector operations.
%
These verifiers, however, encode the
semantics of user-defined functions
with universally-quantified axioms
and use incomplete (albeit effective)
heuristics to instantiate those axioms.
%
These heuristics make it difficult to
characterize the kinds of proofs that
can be automated, and hence,
explain why a given proof attempt
fails~\cite{Leino16}.
%
At the other end, we have
Type-Theory (TT) based theorem provers
(\eg \coq and \agda)
that use type-level computation
(normalization) to facilitate
principled reasoning about
terminating user-defined functions,
but which
% comes at the cost of
require the user to supply
lemmas or rewrite hints
% that add friction to the ubiquitous task of
to discharge proofs
over decidable theories.

%% \RN{As a reader, I don't have that
%%     firm a notion of what the claimed
%%     completeness is at this point.}

We introduce \emph{Refinement Reflection},
a new framework for building SMT-based
deductive verifiers, which permits the
specification of arbitrary properties
and yet enables complete, automated
SMT-based reasoning about user-defined
functions.
%
In previous work, refinement types
\citep{ConstableS87,Rushby98} --- which decorate
basic types (\eg "Integer") with SMT-decidable
predicates (\eg "{v:Integer | 0 <= v && v < 100}") --- were
used to retrofit so-called shallow
verification, such as array bounds checking, into
several languages:
%
ML~\cite{pfenningxi98,GordonRefinement09,LiquidPLDI08},
C~\cite{deputy,LiquidPOPL10},
Haskell~\citep{Vazou14},
TypeScript~\cite{Vekris16}, and
Racket~\cite{RefinedRacket}.
%
In this work, 
we extend refinement types with refinement reflection,
leading to the following three contributions.

\mypara{1. Refinement Reflection}
%
Our first contribution is the notion of
\emph{refinement reflection}.
%
To reason about user-defined functions,
the function's implementation is
\emph{reflected} into its (output)
refinement-type specification,
thus converting the function's
type signature into a precise
description of the function's behavior.
%
This simple idea has a profound consequence:
at \emph{uses} of the function, the standard
rule for (dependent) function application
yields a precise means of reasoning about
the function~(\S~\ref{sec:types-reflection}).


\mypara{2. Complete Specification}
%
% Reflection lets us represent expressive specifications,
% that fall well outside of SMT-decidable logics,
% simply as unit-values refined by logical
% propositions.
%
Our second contribution is a
\emph{library of combinators}
that lets programmers compose
sophisticated \emph{proofs}
from basic refinements and
function definitions.
%
Our proof combinators let programmers
use existing language mechanisms, like
branches (to encode case splits),
recursion (to encode induction), and
functions (to encode auxiliary lemmas),
to write proofs that look very much like
% transcriptions of
their pencil-and-paper
analogues~(\S~\ref{sec:overview}).
%
Furthermore, since proofs are literally
just programs, we use the principle of
propositions-as-types~\cite{Wadler15} 
(known as the Curry-Howard isomorphism~\cite{CHIso})
to show
% , via a recipe for encoding proofs with
% alternating quantifiers,
that SMT-based verifiers can
express any natural deduction proof, 
thus providing a pleasant
implementation of
natural deduction
that can be used for
pedagogical purposes~(\S~\ref{sec:higher-order}).

\mypara{3. Complete Verification}
%
While equational proofs can be very \elegant
and expressive, writing them out can
quickly get exhausting.
%
Our third contribution is
\emph{Proof by Logical Evaluation} (\pbesym)
a new proof-search algorithm
that automates
%that completely automates
equational reasoning.
%
The key idea in \pbesym is to mimic
type-level computation within SMT-logics
by representing functions in a
\emph{guarded form} \cite{Dijkstra1975}
and repeatedly unfolding function application
terms by instantiating them with their definition
corresponding to an \emph{enabled} guard.
%
We formalize a notion of equational proof
and show that the above strategy is
\emph{complete}: \ie it is guaranteed
to find an equational proof if one exists.
%
Furthermore, using techniques from
the literature on Abstract
Interpretation~\cite{CousotCousot77}
and Model Checking \cite{CGL92},
we show that the above proof search
corresponds to a \emph{universal}
(or \emph{must}) abstraction of
the concrete semantics of the
user-defined functions.
%
Thus, since those functions are total,
% and terminating,
we obtain the pleasing
guarantee that proof search
terminates~(\S~\ref{sec:pbe}).

\smallskip
% \mypara{4. Evaluation}
%
We evaluate our approach by implementing
refinement reflection and \pbesym in
\toolname~\citep{Vazou14}, thereby
turning Haskell into a theorem prover.
%
Repurposing an existing programming language allows us
to take advantage of a mature compiler and an ecosystem of
libraries, while keeping proofs and programs in the same
language.
%
We demonstrate the benefits of this
conversion by proving typeclass laws.
%
Haskell's typeclass machinery has led
to a suite of expressive abstractions
and optimizations which, for correctness,
crucially require typeclass \emph{instances}
to obey key algebraic laws.
%
We show how reflection and \pbesym
can be used to verify that widely
used instances of the Monoid,
Applicative, Functor, and Monad
typeclasses satisfy the respective
laws.
%
Finally, we use reflection to
create the first deterministic
parallelism library that actually
verifies assumptions about
associativity and ordering
that ensure determinism (\S~\ref{sec:evaluation}).

Thus, our results demonstrate
that Refinement Reflection and
Proof by Logical Evaluation
identify a new design for
deductive verifiers which,
by combining the complementary
strengths of SMT- and TT- based
approaches, enables complete
verification of expressive
specifications spanning
decidable theories and
user defined functions.
